% =============================================================================
% AGENT   : Geneticist (Domain Expert) — PRIMARY AUTHOR OF THIS MODULE
% MODULE  : Discussion — English
% NOTE    : This module interprets results from the Statistician's outputs.
%           Update whenever statistics are recalculated (Lead Scientist to verify).
% =============================================================================

\section{Discussion}
\label{sec:discussion}

The results obtained from the simulated marker panel align well with
established expectations from the molecular marker literature. The
consistently higher \pic{} and \he{} values observed for \ssr{} markers
corroborate their recognised superiority in capturing polymorphism, a
characteristic attributable to the stepwise mutation model governing
microsatellite repeat evolution and the consequent accumulation of
multiple alleles per locus \autocite{Ellegren2004}.

\snp{} markers, while individually less informative due to their
biallelic nature, compensate through sheer abundance across the
genome. High-density \snp{} genotyping platforms can capture linkage
disequilibrium patterns at fine resolution, enabling genome-wide
association studies (GWAS) and genomic selection with accuracy
surpassing that achievable with sparser \ssr{} panels
\autocite{Goddard2007}. The trade-off between per-locus informativeness
(\ssr{}) and marker density (\snp{}) should therefore be evaluated in
the context of the specific application: parentage analysis and
cultivar identification typically favour \ssr{}, while genomic
prediction and QTL mapping at scale favour dense \snp{} arrays.

The narrow distribution of $H_o$ values around $H_e$ in the simulated
dataset indicates a population in approximate Hardy--Weinberg equilibrium,
as expected given the absence of selection, drift, or inbreeding effects
in the simulation model. In real populations, departures from this
equilibrium may signal biological phenomena of interest, such as
inbreeding depression, selection sweeps, or genotyping artefacts
(e.g., null alleles in \ssr{} loci), which should be systematically
tested prior to downstream analyses \autocite{Peakall2006}.

Taken together, these findings reinforce the complementarity of \snp{}
and \ssr{} marker systems. The modular analytical pipeline developed
here, combining Python-based statistical computation with \LaTeX{}
document generation, provides a reproducible and extensible framework
for future marker characterisation studies.

\subsection{Limitations and Future Directions}

The present analysis is constrained by the use of simulated data and
a biallelic model for \ssr{} loci, which underestimates their true
polymorphism. Future work should incorporate empirical genotyping data
from actual populations, employ multiallelic \pic{} formulations for
\ssr{} markers \autocite{Anderson1993}, and extend the statistical
framework to include $F$-statistics and population structure analyses
(\textit{e.g.}, STRUCTURE, ADMIXTURE).
